\documentclass[letterpaper, 10 pt, conference]{ieeeconf}
\IEEEoverridecommandlockouts
\overrideIEEEmargins

\usepackage[utf8]{inputenc}
\usepackage[T1]{fontenc}
\usepackage{amsmath}
\usepackage{amssymb}
\usepackage{booktabs}
\usepackage{graphicx}
\usepackage{url}

\title{\LARGE \bf
ML4 - Kompresija ML modela\\
}

\author{
  \parbox{4in}{\centering
    Luka Capan$^{1}$, Borna Budić$^{1}$, Silvia Petrović$^{1}$, Jakov Dragičević$^{1}$
    \thanks{$^{1}$Tehnički fakultet, Sveučilište u Rijeci, Hrvatska}
  }
}

\begin{document}

\maketitle
\thispagestyle{empty}
\pagestyle{empty}

%%%%%%%%%%%%%%%%%%%%%%%%%%%%%%%%%%%%%%%%%%%%%%%%%%%%%%%%%%%%%%%%%%%%%%%%%%%%%%%%
\begin{abstract}
Duboke neuronske mreže (DNN) postižu izvanredne rezultate u zadacima klasifikacije slika, no njihova rastuća složenost otežava primjenu na uređajima s ograničenim resursima.
Ovaj rad prikazuje praktičnu implementaciju cjelovite kompresijske cjevovode primijenjene na konvolucijsku neuronsku mrežu (CNN) treniranu od nule na skupu podataka Fruits-360 veličine 100x100 piksela.
Cjevovod kombinira tri komplementarne tehnike: poboljšanu arhitekturu s BatchNorm i Dropout slojevima, dinamički pruning težina tijekom treniranja, te post-training kvantizaciju u INT8.
Cijela implementacija ostvarena je isključivo u biblioteci NumPy, bez korištenja okvira za duboko učenje.
Rezultati pokazuju da druga arhitektura postiže 77.55\% top-1 točnosti na 261 klasi, što je poboljšanje od 5.25 postotnih bodova u odnosu na originalnu osnovicu (72.30\%) uz 52\% manje parametara.
Primjena dinamičkog pruninga zadržava točnost od 77.45\% pri 50\% sparsity, a naknadna INT8 kvantizacija dodatno smanjuje teorijsku veličinu modela 4$\times$ uz pad točnosti od samo 0.05\%.
Konačni komprimirani model je 16.7$\times$ manji od originalne osnovice uz istovremeno bolju točnost.
\end{abstract}

%%%%%%%%%%%%%%%%%%%%%%%%%%%%%%%%%%%%%%%%%%%%%%%%%%%%%%%%%%%%%%%%%%%%%%%%%%%%%%%%
\section{\textbf{UVOD}}

Brzi napredak dubokog učenja proizveo je sve moćnije modele neuronskih mreža, no taj napredak dolazi po cijeni: moderne DNN-ove zahtijevaju značajne memorijske i računalne resurse koji često nisu dostupni na rubnim uređajima poput pametnih telefona, mikrokontrolera i FPGA-a \cite{li2023survey}.
Primjerice, ResNet-50 zahtijeva 98\,MB pohrane i više od 3.8 milijarde operacija u pomičnom zarezu za obradu jedne slike.

Kompresija modela rješava ovaj problem smanjenjem veličine i računalne složenosti trenirane mreže uz zadržavanje što više predikcijske točnosti.
Pregledni rad Li, Li i Meng \cite{li2023survey} identificira pet glavnih obitelji tehnika kompresije: pruning težina, kvantizacija parametara, dekompozicija niskog ranga, distilacija znanja i dizajn lakih modela.

Cilj ovog rada je implementirati i evaluirati cjelovitu kompresijsku cjevovodu koja kombinira tri komplementarne tehnike: \textit{(i)} arhitektonski dizajn s ugrađenom regularizacijom (Sekcija 4 pregleda), \textit{(ii)} dinamički pruning težina (Sekcija 2.2 pregleda), te \textit{(iii)} post-training kvantizaciju (Sekcija 3 pregleda).
Implementacija koristi isključivo NumPy, čime svaka matematička operacija postaje eksplicitna i izravno povezana s jednadžbama iz preglednog rada.

Glavni doprinosi ovog rada su:
\begin{itemize}
  \item Cjelovita implementacija CNN cjevovoda za treniranje u čistom NumPy-u, uključujući propagaciju unaprijed, propagaciju unatrag i SGD s momentumom.
  \item Poboljšana V2 arhitektura s BatchNorm i Dropout slojevima koja smanjuje broj parametara za 52\%.
  \item Dinamički globalni nestrukturirani pruning težina po L1 magnitudi koji postupno povećava sparsity tijekom treninga.
  \item Post-training INT8 simetrična i asimetrična kvantizacija prema Algoritmu 1 i jednadžbama (1)--(5) iz \cite{li2023survey}.
  \item Kvantitativna usporedba točnosti, veličine modela i vremena inferencije kroz cijelu kompresijsku cjevovodu.
\end{itemize}

%%%%%%%%%%%%%%%%%%%%%%%%%%%%%%%%%%%%%%%%%%%%%%%%%%%%%%%%%%%%%%%%%%%%%%%%%%%%%%%%
\section{\textbf{METODOLOGIJA}}

\subsection{Arhitektura mreže — originalna osnovica}

Polazna arhitektura, nazvana \textit{SimpleCNN}, sastoji se od dva konvolucijska bloka iza kojih slijede dva potpuno povezana sloja:

\begin{enumerate}
  \item \textbf{Conv1}: 8 filtara, jezgra $3\times3$, padding 1 $\rightarrow$ ReLU $\rightarrow$ MaxPool $2\times2$
  \item \textbf{Conv2}: 16 filtara, jezgra $3\times3$, padding 1 $\rightarrow$ ReLU $\rightarrow$ MaxPool $2\times2$
  \item \textbf{Flatten}: $(N,\;16,\;25,\;25) \rightarrow (N,\;10000)$
  \item \textbf{FC1}: $10000 \rightarrow 128$, ReLU
  \item \textbf{FC2}: $128 \rightarrow C$ (broj klasa), Softmax
\end{enumerate}

Ulazne slike su RGB rezolucije $100\times100$, normalizirane na $[0,1]$.
Ukupan broj parametara iznosi približno 1.32 milijuna, od čega 1.28 milijuna otpada na sloj FC1.
He inicijalizacija primijenjena je na sve konvolucijske i potpuno povezane slojeve.

\subsection{Poboljšana arhitektura}

Analiza originalnog modela pokazala je značajan overfitting: točnost na treningu doseže približno 99\% dok točnost na testu plateaua oko 72\%.
Razlika od približno 27 postotnih bodova ukazuje na to da model previše ovisi o memorizaciji trening skupa.
Kako bi se ovaj problem riješio, dizajnirana je poboljšana arhitektura s tri konkretne izmjene:

\begin{enumerate}
  \item \textbf{BatchNorm sloj nakon svake konvolucije} -- normalizira aktivacije unutar batcha, stabilizira treniranje i djeluje kao blaga regularizacija \cite{ioffe2015batchnorm}.
  \item \textbf{Dropout sa stopom $p=0.3$ nakon FC1} -- tijekom treniranja slučajno postavlja 30\% aktivacija na nulu, sprječavajući FC1 da memorira pojedinačne primjere.
  \item \textbf{Treći konvolucijski blok (16$\rightarrow$32 filtera)} -- dodaje dubinu i, kao posljedicu, treći MaxPool sloj koji smanjuje prostornu dimenziju s $25\times25$ na $12\times12$. Time se ulaz u FC1 reducira s 10 000 na 4 608 značajki, a sam FC1 sa 1.28M na 590K parametara.
\end{enumerate}

Poboljšana mreža ima ukupno 629 765 parametara, što je smanjenje od 52\% u odnosu na originalnu.
Javno sučelje drugog modela identično je prvom sučelju (forward, backward, get\_trainable\_layers), čime se osigurava da treninzi, pruning i kvantizacija rade bez izmjena ostatka koda.

\subsection{Dinamički pruning težina}

Implementiran je globalni nestrukturirani L1 pruning prema Han et al., kako je opisano u Sekciji 2.2 \cite{li2023survey}.
Za razliku od klasične paradigme \textit{train $\rightarrow$ prune $\rightarrow$ fine-tune}, ovdje primijenjeni \textbf{dinamički pruning} postupno povećava sparsity \textit{tijekom} treniranja prema rasporedu:

\begin{equation}
  s(t) = s_{\text{final}} \cdot \left( 1 - \left(1 - \frac{t - t_0}{t_1 - t_0}\right)^3 \right)
\end{equation}

gdje je $s(t)$ trenutna razina sparsity u epohi $t$, $s_{\text{final}}=0.5$, $t_0=2$ (početna epoha), te $t_1=15$ (završna epoha).
Kubična krivulja osigurava brži rast sparsity u prvim epohama, a usporavanje pri kraju.

Nakon svakog ažuriranja gradijentom, binarna maska se ponovno primjenjuje kako bi se osiguralo da pruned težine ostaju točno nula.
Postupni pristup dopušta modelu da se prilagodi gubitku težina umjesto da iznenadno izgubi 50\% kapaciteta.

\subsection{Post-training kvantizacija}

Implementirana je simetrična i asimetrična INT8 kvantizacija prema Algoritmu 1 i jednadžbama (1)--(5) \cite{li2023survey}.
Za svaki sloj zasebno računaju se faktor skale $S$ i točka pomaka $Z$ iz raspona težina $[R_{\min}, R_{\max}]$:

\begin{equation}
  S = \frac{R_{\max} - R_{\min}}{Q_{\max} - Q_{\min}}, \quad
  Z = Q_{\max} - \frac{R_{\max}}{S}
\end{equation}

Težine se zatim kvantiziraju u INT8:

\begin{equation}
  Q = \mathrm{round}\!\left(\frac{R}{S} + Z\right), \quad Q \in [-128,\,127]
\end{equation}

i dekvantiziraju natrag u FP32 za inferenciju:

\begin{equation}
  \hat{R} = (Q - Z) \cdot S
\end{equation}

Round-trip postupak FP32$\rightarrow$INT8$\rightarrow$FP32 omogućuje mjerenje utjecaja kvantizacije na točnost dok forward prolaz i dalje koristi FP32 aritmetiku.
U stvarnoj hardverskoj implementaciji koja sprema težine kao INT8 (1 bajt umjesto 4), ovo daje teoretsko 4$\times$ smanjenje veličine modela.

%%%%%%%%%%%%%%%%%%%%%%%%%%%%%%%%%%%%%%%%%%%%%%%%%%%%%%%%%%%%%%%%%%%%%%%%%%%%%%%%
\section{\textbf{EKSPERIMENTALNI POSTAV}}

\subsection{Skup podataka}

Korišten je \textit{Fruits-360} \cite{fruits360}, javno dostupan skup slika voća snimljenih nad bijelom pozadinom.
Svaka slika je rezolucije $100\times100$ piksela u RGB formatu.
Korištena verzija sadrži 261 različitu klasu voća i povrća.

Za eksperimente korišteno je 100 slika po klasi za treniranje i potpuni test skup za evaluaciju (približno 26 016 trening i 25 244 test slika).
Slike su normalizirane na $[0,1]$ dijeljenjem vrijednosti piksela s 255.
Nije primijenjena augmentacija podataka.

\subsection{Tijek eksperimenta}

Cjelovita kompresijska cjevovoda izvodi se u tri koraka:
\begin{enumerate}
  \item \textbf{Treniranje s dinamičkim pruningom}: SimpleCNNv2 trenira se 15 epoha s $\eta = 0.01$, batch size 32, momentum 0.9. Pruning počinje u 2. epohi i postupno postiže 50\% sparsity do 15. epohe.
  \item \textbf{Kvantizacija}: učitava se najbolji checkpoint iz koraka 1, a zatim se primjenjuje asimetrična INT8 PTQ sloj po sloj.
  \item \textbf{Evaluacija}: točnost se mjeri na potpunom test skupu nakon svakog koraka.
\end{enumerate}

Cjevovod je pokrenut na laptopu s 16\,GB RAM-a bez GPU akceleracije. Ukupno trajanje pravog runa iznosi približno 5 sati i 30 minuta.

%%%%%%%%%%%%%%%%%%%%%%%%%%%%%%%%%%%%%%%%%%%%%%%%%%%%%%%%%%%%%%%%%%%%%%%%%%%%%%%%
\section{\textbf{REZULTATI I RASPRAVA}}

\subsection{Usporedba arhitektura modela}

Tablica~\ref{tab:v1_vs_v2} prikazuje rezultate usporedbe modela arhitektura pod identičnim uvjetima treniranja (isti seed, isti dataset, 10 epoha, 100 slika po klasi).

\begin{table}[h]
\caption{Usporedba modela arhitekture pod identičnim uvjetima.}
\label{tab:v1_vs_v2}
\begin{center}
\begin{tabular}{lcc}
\toprule
\textbf{Metrika} & \textbf{Original} & \textbf{Poboljšana} \\
\midrule
Top-1 točnost (\%)     & 72.30  & \textbf{77.55} \\
Broj parametara        & 1 315 189 & \textbf{629 765} \\
Veličina modela (KB)   & 10 275  & 4 920 \\
Vrijeme inferencije (ms) & 82.30 & 89.14 \\
Vrijeme treninga       & 2h 34min & 3h 16min \\
\midrule
Promjena točnosti      & --- & \textbf{+5.25\%} \\
Promjena broja parametara & --- & \textbf{$-$52.1\%} \\
\bottomrule
\end{tabular}
\end{center}
\end{table}

Poboljšana arhitektura postiže poboljšanje od 5.25 postotnih bodova u točnosti uz istovremeno smanjenje broja parametara za 52.1\%.
Iako je inferencija 8.3\% sporija zbog BatchNorm operacija, a trening 27\% duži zbog dodatnog konvolucijskog sloja, oba ova troška su jednokratna i prihvatljiva za poboljšanje generalizacije.

\subsection{Dinamički pruning na poboljšanom modelu}

Tablica~\ref{tab:pruning_progress} prikazuje napredak treninga s dinamičkim pruningom kroz 15 epoha.

\begin{table}[h]
\caption{Napredak dinamičkog pruninga tijekom treniranja V2 modela.}
\label{tab:pruning_progress}
\begin{center}
\begin{tabular}{cccc}
\toprule
\textbf{Epoha} & \textbf{Gubitak} & \textbf{Točnost (\%)} & \textbf{Sparsity (\%)} \\
\midrule
1   & 5.04 &  9.39 & 0.00 \\
3   & 1.95 & 64.25 & 18.44 \\
5   & 0.94 & 71.73 & 31.75 \\
7   & 0.60 & 74.31 & 40.67 \\
10  & 0.36 & 75.89 & 47.72 \\
12  & 0.27 & 76.81 & 49.50 \\
14  & 0.23 & \textbf{77.45} & 49.98 \\
15  & 0.20 & 77.26 & 50.00 \\
\bottomrule
\end{tabular}
\end{center}
\end{table}

Najbolji model postignut je u 14. epohi s 77.45\% točnosti pri sparsity od 49.98\%.
Bitno je primijetiti da pruning praktički nije naštetio točnosti: poboljšani bez pruninga postiže 77.55\%, a poboljšani s 50\% pruningom 77.45\% --- razlika od samo 0.10 postotnih bodova.
Postupna priroda dinamičkog pruninga omogućila je modelu da se prilagodi gubitku težina umjesto da podnese iznenadnu redukciju kapaciteta.

\subsection{Konačna kompresija s kvantizacijom}

Nakon kvantizacije pruned poboljšanog modela u INT8, dobiveni su rezultati prikazani u Tablici~\ref{tab:quant}.

\begin{table}[h]
\caption{Rezultati INT8 kvantizacije primijenjene na pruned poboljšani model.}
\label{tab:quant}
\begin{center}
\begin{tabular}{lcc}
\toprule
\textbf{Metrika} & \textbf{Pruned V2 (FP32)} & \textbf{Quantizirano (INT8)} \\
\midrule
Točnost (\%)         & 76.53 & 76.48 \\
Veličina (KB)         & 4 920 & \textbf{614.5} (teor.) \\
Vrijeme inferencije (ms) & 87.36 & \textbf{73.60} \\
\midrule
Pad točnosti          & --- & 0.05\% \\
Ubrzanje              & --- & 1.19$\times$ \\
\bottomrule
\end{tabular}
\end{center}
\end{table}

INT8 kvantizacija postiže teorijsko 4$\times$ smanjenje veličine modela uz zanemariv pad točnosti od 0.05\%.
Mjereno ubrzanje inferencije od 1.19$\times$ posljedica je smanjenog dinamičkog raspona težina nakon kvantizacije.

\subsection{Cjelovita usporedba: Modeli → Pruning → Kvantizacija}

Tablica~\ref{tab:full_pipeline} prikazuje konačnu usporedbu svih faza kroz cijelu kompresijsku cjevovodu.

\begin{table}[h]
\caption{Cjelovita usporedba svih faza kompresije u odnosu na V1 (originalnu) osnovicu.}
\label{tab:full_pipeline}
\begin{center}
\footnotesize
\begin{tabular}{lccc}
\toprule
\textbf{Faza} & \textbf{Točnost (\%)} & \textbf{Veličina (KB)} & \textbf{Faktor} \\
\midrule
V1 osnovica (KT1)            & 72.30 & 10 275 & 1.0$\times$ \\
V2 (BN + Dropout)            & 77.55 & 4 920  & 2.1$\times$ \\
V2 + Pruning                 & 77.45 & 2 460$^{*}$ & 4.2$\times$ \\
V2 + Pruning + INT8          & \textbf{76.48} & \textbf{614.5} & \textbf{16.7$\times$} \\
\bottomrule
\end{tabular}
\\[2pt]
{\scriptsize $^{*}$Efektivna veličina uz 50\% nula u težinama.}
\end{center}
\end{table}

Konačni komprimirani model je \textbf{16.7$\times$ manji} od originalne V1 osnovice, dok istovremeno postiže \textbf{4.18 postotnih bodova bolju točnost} (76.48\% vs 72.30\%).
Ovaj rezultat pokazuje da dobro projektirana kombinacija arhitektonskog dizajna i tehnika kompresije može istovremeno poboljšati i sažeti model, što je u skladu s preporukama iz pregleda \cite{li2023survey}.

\subsection{Per-layer analiza kvantizacijske greške}

Prosječna greška kvantizacije po sloju (srednja apsolutna razlika između originalnih i dekvantiziranih težina) ostala je ispod $0.003$ za sve slojeve.
Najmanja greška opažena je u sloju FC1 ($0.000377$), što potvrđuje da je 256-razinska INT8 rešetka više nego dovoljno fina za distribucije težina opažene u potpuno povezanim slojevima.

%%%%%%%%%%%%%%%%%%%%%%%%%%%%%%%%%%%%%%%%%%%%%%%%%%%%%%%%%%%%%%%%%%%%%%%%%%%%%%%%
\section{\textbf{ZAKLJUČAK}}

Ovaj rad prikazao je cjelovitu kompresijsku cjevovodu za CNN modele implementiranu u čistom NumPy-u, kombinirajući tri komplementarne tehnike opisane u pregledu Li, Li i Meng \cite{li2023survey}.

Glavni rezultati su sljedeći.
Poboljšana arhitektura s BatchNorm i Dropout slojevima postigla je 77.55\% top-1 točnosti uz 52\% manje parametara od originalne osnovice.
Dinamički L1 pruning sa 50\% sparsity zadržao je točnost na 77.45\% (pad od samo 0.10 postotnih bodova).
Post-training INT8 kvantizacija ostvarila je teorijsko 4$\times$ smanjenje veličine uz pad točnosti od 0.05\% i ubrzanje inferencije od 1.19$\times$.

Kombinacija sve tri tehnike daje finalni model koji je 16.7$\times$ manji od originalne osnovice uz istovremeno 4.18 postotnih bodova \textit{bolju} top-1 točnost.
Time je pokazano da pažljivim arhitektonskim dizajnom u kombinaciji s dinamičkim pruningom i INT8 kvantizacijom moguće postići značajnu kompresiju bez ikakvog značajnog gubitka prediktivne sposobnosti.

%%%%%%%%%%%%%%%%%%%%%%%%%%%%%%%%%%%%%%%%%%%%%%%%%%%%%%%%%%%%%%%%%%%%%%%%%%%%%%%%

\begin{thebibliography}{99}

\bibitem{li2023survey}
Z.\ Li, H.\ Li, and L.\ Meng,
,,Model Compression for Deep Neural Networks: A Survey,''
\textit{Computers}, vol.~12, no.~3, p.~60, ožujak 2023.
\url{https://doi.org/10.3390/computers12030060}

\bibitem{fruits360}
H.\ Mure\c{s}an and M.\ Oltean,
,,Fruit recognition from images using deep learning,''
\textit{Acta Universitatis Sapientiae, Informatica},
vol.~10, br.~1, str.~26--42, 2018.

\bibitem{han2015learning}
S.\ Han, J.\ Pool, J.\ Tran, and W.\ Dally,
,,Learning both weights and connections for efficient neural network,''
u \textit{Proc.\ NeurIPS}, Montreal, QC, Kanada, 2015, str.~1135--1143.

\bibitem{ioffe2015batchnorm}
S.\ Ioffe and C.\ Szegedy,
,,Batch normalization: Accelerating deep network training by reducing internal covariate shift,''
u \textit{Proc.\ ICML}, Lille, Francuska, 2015, str.~448--456.

\bibitem{nagel2021white}
M.\ Nagel, M.\ Fournarakis, R.\ A.\ Amjad, Y.\ Bondarenko,
M.\ van Baalen, and T.\ Blankevoort,
,,A white paper on neural network quantization,''
\textit{arXiv preprint arXiv:2106.08295}, 2021.

\end{thebibliography}

\end{document}
